\documentclass[10pt,a4paper]{amsart}
\usepackage{amsmath,amssymb,amsthm,graphicx}
% For commutative diagrams we use Xy-pic
\usepackage[all]{xy}

% We use 2cell for 2-commutative diagrams.
\xyoption{2cell}
% Theorem environments.
%
\theoremstyle{plain}
\newtheorem{theorem}[subsection]{Theorem}
\newtheorem{proposition}[subsection]{Proposition}
\newtheorem{lemma}[subsection]{Lemma}

\theoremstyle{definition}
\newtheorem{definition}[subsection]{Definition}
\newtheorem{example}[subsection]{Example}
\newtheorem{exercise}[subsection]{Exercise}
\newtheorem{situation}[subsection]{Situation}

\theoremstyle{remark}
\newtheorem{remark}[subsection]{Remark}
\newtheorem{remarks}[subsection]{Remarks}

\numberwithin{equation}{subsection}

% Macros
%
\def\lim{\mathop{\mathrm{lim}}\nolimits}
\def\colim{\mathop{\mathrm{colim}}\nolimits}
\def\Spec{\mathop{\mathrm{Spec}}}
\def\Hom{\mathop{\mathrm{Hom}}\nolimits}
\def\Ext{\mathop{\mathrm{Ext}}\nolimits}
\def\SheafHom{\mathop{\mathcal{H}\!\mathit{om}}\nolimits}
\def\SheafExt{\mathop{\mathcal{E}\!\mathit{xt}}\nolimits}
\def\Sch{\mathit{Sch}}
\def\Mor{\mathop{\mathrm{Mor}}\nolimits}
\def\Ob{\mathop{\mathrm{Ob}}\nolimits}
\def\Sh{\mathop{\mathit{Sh}}\nolimits}
\def\NL{\mathop{N\!L}\nolimits}
\def\CH{\mathop{\mathrm{CH}}\nolimits}
\def\proetale{{pro\text{-}\acute{e}tale}}
\def\etale{{\acute{e}tale}}
\def\QCoh{\mathit{QCoh}}
\def\Ker{\mathop{\mathrm{Ker}}}
\def\Im{\mathop{\mathrm{Im}}}
\def\Coker{\mathop{\mathrm{Coker}}}
\def\Coim{\mathop{\mathrm{Coim}}}

% Boxtimes
%
\DeclareMathSymbol{\boxtimes}{\mathbin}{AMSa}{"02}

%
% Macros for moduli stacks/spaces
%
\def\QCohstack{\mathcal{QC}\!\mathit{oh}}
\def\Cohstack{\mathcal{C}\!\mathit{oh}}
\def\Spacesstack{\mathcal{S}\!\mathit{paces}}
\def\Quotfunctor{\mathrm{Quot}}
\def\Hilbfunctor{\mathrm{Hilb}}
\def\Curvesstack{\mathcal{C}\!\mathit{urves}}
\def\Polarizedstack{\mathcal{P}\!\mathit{olarized}}
\def\Complexesstack{\mathcal{C}\!\mathit{omplexes}}
% \Pic is the operator that assigns to X its picard group, usage \Pic(X)
% \Picardstack_{X/B} denotes the Picard stack of X over B
% \Picardfunctor_{X/B} denotes the Picard functor of X over B
\def\Pic{\mathop{\mathrm{Pic}}\nolimits}
\def\Picardstack{\mathcal{P}\!\mathit{ic}}
\def\Picardfunctor{\mathrm{Pic}}
\def\Deformationcategory{\mathcal{D}\!\mathit{ef}}
\usepackage{cleveref}
\crefname{originalreference}{original reference}{original references}
\Crefname{originalreference}{Original reference}{Original references}
\makeatletter
\newlabel{lemma-proper-base-change-f-star}{{0A3U}{source}{}{sourcepointer-1}{}}
\newlabel{lemma-proper-base-change-f-star@cref}{{[originalreference][0][]0A3U}{[1][0][]source}{}{}{}}
\makeatother
\crefname{originalreference}{original reference}{original references}
\Crefname{originalreference}{Original reference}{Original references}
\makeatletter
\newlabel{lemma-proper-base-change-in-terms-of-injectives}{{0A4B}{source}{}{sourcepointer-2}{}}
\newlabel{lemma-proper-base-change-in-terms-of-injectives@cref}{{[originalreference][0][]0A4B}{[1][0][]source}{}{}{}}
\makeatother
\crefname{originallemma}{original lemma}{original lemmas}
\Crefname{originallemma}{Original lemma}{Original lemmas}
\makeatletter
\newlabel{sites-cohomology-lemma-relative-Leray}{{0734}{source}{}{sourcepointer-3}{}}
\newlabel{sites-cohomology-lemma-relative-Leray@cref}{{[originallemma][0][]0734}{[1][0][]source}{}{}{}}
\makeatother
\title{Base change for composites: if cohomology commutes with base change for f and for g o f with f surjective, then it does for g}
\author{The Stacks Project Authors}
\date{stacks-project:etale-cohomology.tex:local-snapshot-20261011; GNU FDL 1.2 or later}
\begin{document}
\maketitle
\noindent\textit{Editorial note.} All delivered mathematical text is the author's own, copied from the pinned Stacks Project chapter file etale-cohomology.tex; no statement or proof step was rewritten, reordered or removed.\par\medskip
\noindent\textit{Editorial note.} The counted proof is the author's complete argument: the proof is written out in full: the reduction from a general injective sheaf to f\_*J using surjectivity of f, the base change square with the pushforward base change formula, and the degeneration of the relative Leray spectral sequence under the two vanishing hypotheses.\par\medskip
\noindent\textit{Editorial note.} The label lemma-sandwich names the counted statement, and the proof environment immediately following it in the source is delivered verbatim as its proof.\par\medskip
\noindent\textit{Editorial note.} The counted statement is carried by the author's own lemma environment in the chapter source.\par\medskip
\noindent\textit{Editorial note.} The statement and proof are copied from the pinned chapter file; the three cited inputs are kept as external pointers with their canonical Stacks tags.\par\medskip
\noindent\textit{Editorial note.} The preamble is the chapter's own native preamble excerpt already used for this source: the Xy-pic setup (native preamble lines 14-18) and the native theorem declarations with the mathematical macros (native preamble lines 154-217).\par\medskip
\noindent\textit{Allowed standard prerequisites (editorial).}\begin{itemize}
\item Base change in terms of injectives for torsion sheaves. Same chapter, the reformulation of the base change property through injective sheaves of Z/ellZ-modules. Supplies the form of the property used throughout: it suffices that the higher direct images of the relevant injective sheaves vanish after base change.
\item Base change for the pushforward along a proper morphism. Same chapter, the lemma b\textasciicircum{}\{-1\}f\_*J = f'\_*a\textasciicircum{}\{-1\}J for the base change square of f. Supplies the identification used in the second half of the proof.
\item The relative Leray spectral sequence. Cohomology on Sites, the relative Leray spectral sequence for a composition of morphisms. Supplies the spectral sequence R\textasciicircum{}p g'\_* R\textasciicircum{}q f'\_* -> R\textasciicircum{}\{p+q\}(g' o f')\_* whose vanishing input gives the conclusion.
\end{itemize}\medskip
\ifdefined\hypertarget\hypertarget{sourcepointer-1}{}\fi
\noindent\textit{Source reference.} Original reference 0A3U: sources/stacks-etale-cohomology/pinned-original-source.tex, source lines 18282--18282. Prerequisite pointer invoked at proof line(s) 18461; it names a standard fact of the same chapter that the counted argument uses, and every step of the counted proof built on it is retained verbatim in the delivered body.\par\medskip
\ifdefined\hypertarget\hypertarget{sourcepointer-2}{}\fi
\noindent\textit{Source reference.} Original reference 0A4B: sources/stacks-etale-cohomology/pinned-original-source.tex, source lines 18360--18360. Prerequisite pointer invoked at proof line(s) 18443; it names a standard fact of the same chapter that the counted argument uses, and every step of the counted proof built on it is retained verbatim in the delivered body.\par\medskip
\ifdefined\hypertarget\hypertarget{sourcepointer-3}{}\fi
\noindent\textit{Source reference.} Original lemma 0734: sources/stacks-chow/tags.txt, source lines 9048--9048. Prerequisite pointer invoked at proof line(s) 18466; it names a standard fact of a sibling chapter that the counted argument uses, and every step of the counted proof built on it is retained verbatim in the delivered body.\par\medskip
\begin{lemma}
\label{lemma-sandwich}
Let $f : X \to Y$ and $g : Y \to Z$ be proper morphisms of schemes. Assume
\begin{enumerate}
\item cohomology commutes with base change for $f$,
\item cohomology commutes with base change for $g \circ f$, and
\item $f$ is surjective.
\end{enumerate}
Then cohomology commutes with base change for $g$.
\end{lemma}

\begin{proof}
We will use the equivalence of
Lemma \ref{lemma-proper-base-change-in-terms-of-injectives}
without further mention. Let $\ell$ be a prime number.
Let $\mathcal{I}$ be an injective sheaf of
$\mathbf{Z}/\ell\mathbf{Z}$-modules on $Y_\etale$.
Choose an injective map of sheaves $f^{-1}\mathcal{I} \to \mathcal{J}$
where $\mathcal{J}$ is an injective sheaf of
$\mathbf{Z}/\ell\mathbf{Z}$-modules on $X_\etale$.
Since $f$ is surjective the map $\mathcal{I} \to f_*\mathcal{J}$
is injective (look at stalks in geometric points).
Since $\mathcal{I}$ is injective we see that $\mathcal{I}$
is a direct summand of $f_*\mathcal{J}$. Thus it suffices
to prove the desired vanishing for $f_*\mathcal{J}$.

\medskip\noindent
Let $Z' \to Z$ be a morphism of schemes and set
$Y' = Z' \times_Z Y$ and $X' = Z' \times_Z X = Y' \times_ Y X$.
Denote $a : X' \to X$, $b : Y' \to Y$, and $c : Z' \to Z$ the
projections. Similarly for $f' : X' \to Y'$ and $g' : Y' \to Z'$.
By Lemma \ref{lemma-proper-base-change-f-star} we have
$b^{-1}f_*\mathcal{J} = f'_*a^{-1}\mathcal{J}$.
On the other hand, we know that $R^qf'_*a^{-1}\mathcal{J}$ and
$R^q(g' \circ f')_*a^{-1}\mathcal{J}$ are zero for $q > 0$.
Using the spectral sequence
(Cohomology on Sites, Lemma \ref{sites-cohomology-lemma-relative-Leray})
$$
R^pg'_* R^qf'_* a^{-1}\mathcal{J} \Rightarrow
R^{p + q}(g' \circ f')_* a^{-1}\mathcal{J}
$$
we conclude that
$ R^pg'_*(b^{-1}f_*\mathcal{J}) = R^pg'_*(f'_*a^{-1}\mathcal{J}) = 0$
for $p > 0$ as desired.
\end{proof}

\end{document}
